% archon:covers AlgebraicJacobian/RiemannRoch/CurveChartDedekind.lean

\chapter{An affine chart of a smooth curve is a Dedekind domain}
\label{chap:RiemannRoch_CurveChartDedekind}

\providecommand{\Spec}{\operatorname{Spec}}
\providecommand{\krullDim}{\dim}
\providecommand{\coheight}{\operatorname{coheight}}
\providecommand{\Frac}{\operatorname{Frac}}
\providecommand{\Over}{\operatorname{Over}}

% SOURCE: [Hartshorne], I.6, p. 40 (Theorem 6.2A and the Definition of a
% Dedekind domain) and p. 41 (the nonsingular-curve discussion). Read from
% references/hartshorne-algebraic-geometry.pdf.

\section{Purpose and standing hypotheses}

This chapter isolates the single commutative-algebra crux on which the
structure-sheaf bridge of \cref{chap:RiemannRoch_OcOfD} rests: the coordinate
ring of an affine chart of a smooth proper curve is a \emph{Dedekind domain}.
That fact is what licenses the identification, on each affine open, of the
sections of the zero carrier sheaf with the structure sheaf via the
algebraic-Hartogs principle \cref{lem:mem_integers_of_valuation_le_one} (which
demands a Dedekind base ring). Isolating the crux in its own chapter gives the
bridge a clean \(\texttt{IsDedekindDomain}\) interface and lets
the Dedekind argument be developed independently of the sheaf-gluing work.

\paragraph{Standing hypotheses.} Throughout, \(\bar k\) is an algebraically
closed field and \(C : \Over(\Spec \bar k)\) is a curve over \(\bar k\) whose
structure morphism \(C.\mathrm{hom} : C \to \Spec \bar k\) is smooth of
relative dimension one (\(\texttt{SmoothOfRelativeDimension}\ 1\ C.\mathrm{hom}\))
and locally of finite type (\(\texttt{LocallyOfFiniteType}\ C.\mathrm{hom}\)),
and whose underlying scheme \(C.\mathrm{left}\) is integral
(\(\texttt{IsIntegral}\ C.\mathrm{left}\)). Fix a nonempty affine open
\(U \subseteq C\) (\(\texttt{IsAffineOpen}\ U\)) and write
\[
  A \;=\; \Gamma(C, U)
\]
for its coordinate ring, so that \(U \cong \Spec A\). Because \(C.\mathrm{left}\)
is integral and \(U\) is a nonempty open, \(A\) is an integral domain
(\(\texttt{IsDomain}\ A\)), and its fraction field is the function field of the
curve, \(\Frac(A) = K(C)\) (\(\texttt{IsFractionRing}\ A\ K(C)\)).

\paragraph{Strategy.} We verify the three defining properties of a Dedekind
domain separately and then assemble them through Mathlib's characterisation
\(\texttt{isDedekindRing\_iff}\):
\begin{enumerate}
  \item \emph{Dimension at most one.} The Krull dimension of \(A\) is bounded
    by the Krull dimension of the curve, which is at most one by the project
    bound \cref{thm:krullDim_curve_le_one}; hence every nonzero prime of \(A\)
    is maximal.
  \item \emph{Integral closure.} Integral-closedness is a local property
    (\cref{lem:isIntegrallyClosed_of_localization_maximal_mathlib}); at each
    nonzero maximal ideal the localisation is the curve stalk at a
    codimension-one point, which is a discrete valuation ring
    (\cref{lem:stalk_isDVR_of_smooth}) and therefore integrally closed.
  \item \emph{Noetherianity} comes for free from local Noetherianity of the
    curve, the structure morphism being locally of finite type over the field
    \(\bar k\).
\end{enumerate}

\section{Mathlib dependency anchors}

\begin{theorem}[Dedekind characterisation of a domain]
  \label{lem:isDedekindRing_iff_mathlib}
  \lean{isDedekindRing_iff}
  \mathlibok
  \textit{Provided by Mathlib (\texttt{Mathlib.RingTheory.DedekindDomain.Basic}).}
  Let \(A\) be a commutative ring and \(K\) a fraction field of \(A\)
  (\(\texttt{IsFractionRing}\ A\ K\)). Then \(A\) is a Dedekind ring if and
  only if it is Noetherian, of dimension at most one, and integrally closed in
  \(K\):
  \[
    \texttt{IsDedekindRing}\ A \;\Longleftrightarrow\;
    \texttt{IsNoetherianRing}\ A \;\wedge\;
    \texttt{Ring.DimensionLEOne}\ A \;\wedge\;
    \bigl(\forall\, x \in K,\ \text{$x$ integral over $A$}
      \Rightarrow x \in (\operatorname{algebraMap} A\,K).\mathrm{range}\bigr).
  \]
\end{theorem}

\begin{lemma}[Integral closure is a local property at maximal ideals]
  \label{lem:isIntegrallyClosed_of_localization_maximal_mathlib}
  \lean{IsIntegrallyClosed.of_localization_maximal}
  \mathlibok
  \textit{Provided by Mathlib
    (\texttt{Mathlib.RingTheory.LocalProperties.IntegrallyClosed}).}
  Let \(R\) be an integral domain. Suppose that for every nonzero maximal ideal
  \(\mathfrak p \subseteq R\) the localisation \(\texttt{Localization.AtPrime}\
  \mathfrak p\) is integrally closed. Then \(R\) is integrally closed
  (\(\texttt{IsIntegrallyClosed}\ R\)).
\end{lemma}

\section{Dimension at most one}

\begin{lemma}\leanok
[The chart ring has dimension at most one]
  \label{lem:chart_dimensionLEOne}
  \lean{AlgebraicGeometry.chart_dimensionLEOne}
  \uses{thm:krullDim_curve_le_one}
  Under the standing hypotheses, the coordinate ring \(A = \Gamma(C, U)\) of a
  nonempty affine open of the smooth curve \(C\) has Krull dimension at most
  one: every nonzero prime ideal of \(A\) is maximal
  (\(\texttt{Ring.DimensionLEOne}\ A\)).
\end{lemma}

\begin{proof}
  \uses{thm:krullDim_curve_le_one}
  \leanok
  The affine open \(U\) is homeomorphic to \(\Spec A\), and the inclusion
  \(U \hookrightarrow C.\mathrm{left}\) is an open immersion of topological
  spaces. The Krull dimension of a topological space does not increase under
  passage to a subspace, so
  \[
    \krullDim(\Spec A) \;=\; \krullDim U \;\le\; \krullDim C.\mathrm{left}.
  \]
  By the project bound \cref{thm:krullDim_curve_le_one} the right-hand side is
  at most one, whence \(\krullDim A \le 1\). For a ring this bound is exactly
  the statement that every nonzero prime ideal is maximal: a chain
  \(\mathfrak p_0 \subsetneq \mathfrak p_1 \subsetneq \mathfrak p_2\) of primes
  would witness Krull dimension at least two, so a nonzero prime
  \(\mathfrak p\) (which strictly contains the zero prime, \(A\) being a
  domain) admits no prime strictly above it and is therefore maximal. This is
  \(\texttt{Ring.DimensionLEOne}\ A\).
\end{proof}

\section{Integral closure}

\begin{lemma}\leanok
[The chart ring is integrally closed]
  \label{lem:chart_isIntegrallyClosed}
  \lean{AlgebraicGeometry.chart_isIntegrallyClosed}
  \uses{lem:isIntegrallyClosed_of_localization_maximal_mathlib,
        lem:stalk_isDVR_of_smooth}
  % SOURCE: [Hartshorne], I.6, Theorem 6.2A, p. 40 (read from
  % references/hartshorne-algebraic-geometry.pdf): a noetherian local domain of
  % dimension one is integrally closed iff it is a discrete valuation ring;
  % equivalently a discrete valuation ring is integrally closed.
  % SOURCE QUOTE: "Theorem 6.2A. Let $A$ be a noetherian local domain of
  % dimension one, with maximal ideal $\mathfrak m$. Then the following
  % conditions are equivalent: (i) $A$ is a discrete valuation ring; (ii) $A$
  % is integrally closed; (iii) $A$ is a regular local ring; (iv) $\mathfrak m$
  % is a principal ideal."
  \textit{Source: Hartshorne, I.6, Theorem 6.2A.}
  Under the standing hypotheses, the coordinate ring \(A = \Gamma(C, U)\) is
  integrally closed in its fraction field
  (\(\texttt{IsIntegrallyClosed}\ A\)).
\end{lemma}

\begin{proof}
  \uses{lem:isIntegrallyClosed_of_localization_maximal_mathlib,
        lem:stalk_isDVR_of_smooth}
        \leanok
  Since \(A\) is an integral domain, by
  \cref{lem:isIntegrallyClosed_of_localization_maximal_mathlib} it suffices to
  prove that for every nonzero maximal ideal \(\mathfrak p \subseteq A\) the
  localisation \(A_{\mathfrak p} = \texttt{Localization.AtPrime}\ \mathfrak p\)
  is integrally closed.

  Fix such a \(\mathfrak p\). Under the homeomorphism \(U \cong \Spec A\) the
  maximal ideal \(\mathfrak p\) corresponds to a closed point
  \(x \in U \subseteq C\), and the localisation of the coordinate ring of an
  affine open at the prime of a point is canonically the stalk of the structure
  sheaf there:
  \[
    A_{\mathfrak p} \;\cong\; \mathcal O_{C, x}.
  \]
  Because \(x\) is a closed point of the one-dimensional integral curve \(C\),
  its local ring has Krull dimension one, i.e.\ \(x\) is a codimension-one
  point (\(\coheight x = 1\)). The smooth-curve discrete-valuation witness
  \cref{lem:stalk_isDVR_of_smooth} then makes \(\mathcal O_{C, x}\) a discrete
  valuation ring. A discrete valuation ring is a principal ideal domain, hence
  integrally closed; transporting along the isomorphism above,
  \(A_{\mathfrak p}\) is integrally closed. As \(\mathfrak p\) was an arbitrary
  nonzero maximal ideal, \cref{lem:isIntegrallyClosed_of_localization_maximal_mathlib}
  yields \(\texttt{IsIntegrallyClosed}\ A\).
\end{proof}

\section{The chart ring is a Dedekind domain}

\begin{theorem}\leanok
[An affine chart of a smooth curve is a Dedekind domain]
  \label{thm:chart_isDedekindDomain}
  \lean{AlgebraicGeometry.chart_isDedekindDomain}
  \uses{lem:chart_dimensionLEOne, lem:chart_isIntegrallyClosed,
        lem:isDedekindRing_iff_mathlib}
  % SOURCE: [Hartshorne], I.6, p. 40 (Definition of a Dedekind domain) and
  % p. 41 (nonsingular curves have Dedekind coordinate rings) (read from
  % references/hartshorne-algebraic-geometry.pdf).
  % SOURCE QUOTE: "Definition. A Dedekind domain is an integrally closed
  % noetherian domain of dimension one."
  % SOURCE QUOTE: "If $P$ is a point on a nonsingular curve $Y$, then by (5.1)
  % the local ring $\mathcal{O}_P$ is a regular local ring of dimension one,
  % and so by (6.2A) it is a discrete valuation ring. Its quotient field is the
  % function field $K$ of $Y$, and since $k \subseteq \mathcal{O}_P$, it is a
  % valuation ring of $K/k$."
  \textit{Source: Hartshorne, I.6, Definition (Dedekind domain) and p. 41.}
  Under the standing hypotheses, the coordinate ring \(A = \Gamma(C, U)\) of a
  nonempty affine open of the smooth proper curve \(C\) is a Dedekind domain
  (\(\texttt{IsDedekindDomain}\ A\), equivalently \(\texttt{IsDedekindRing}\ A\)),
  with fraction field the function field \(K(C)\).
\end{theorem}

\begin{proof}
  \uses{lem:chart_dimensionLEOne, lem:chart_isIntegrallyClosed,
        lem:isDedekindRing_iff_mathlib}
        \leanok
  We apply the Dedekind characterisation
  \cref{lem:isDedekindRing_iff_mathlib} to \(A\) with its fraction field
  \(K = K(C)\) (\(\texttt{IsFractionRing}\ A\ K(C)\), since \(K(C)\) is the
  function field of the integral curve and \(A\) is the section ring of a
  nonempty affine open). It remains to verify the three conjuncts.

  \emph{Noetherian.} The structure morphism \(C.\mathrm{hom}\) is locally of
  finite type over the field \(\bar k\), so \(C.\mathrm{left}\) is locally
  Noetherian; the coordinate ring of an affine open of a locally Noetherian
  scheme is a Noetherian ring, giving \(\texttt{IsNoetherianRing}\ A\).

  \emph{Dimension at most one.} This is \cref{lem:chart_dimensionLEOne}:
  \(\texttt{Ring.DimensionLEOne}\ A\).

  \emph{Integrally closed in \(K(C)\).} By
  \cref{lem:chart_isIntegrallyClosed} the ring \(A\) is integrally closed,
  which is exactly the third conjunct: every \(x \in K(C)\) integral over \(A\)
  lies in the image of \(A \hookrightarrow K(C)\).

  The three conjuncts assembled through
  \cref{lem:isDedekindRing_iff_mathlib} give \(\texttt{IsDedekindRing}\ A\);
  together with the domain structure on \(A\) this is
  \(\texttt{IsDedekindDomain}\ A\). The fraction field is \(K(C)\) by
  construction.
\end{proof}

\section{Algebraic Hartogs on a Dedekind chart}

The Dedekind structure of the previous section delivers the converse of the
``regular \(\Rightarrow\) no poles'' inclusion: a rational function that is
\emph{pole-free} on the chart is in fact regular there. We package this
algebraic-Hartogs statement at the level of the coordinate ring, decoupled from
the structure-sheaf gluing of \cref{chap:RiemannRoch_OcOfD}, so that the
sheaf bridge can consume it as a single self-contained fact. It is the affine,
codimension-one shadow of the principle ``intersection of localisations \(=\)
the ring'' (Hartshorne~II.6.3A).

\begin{lemma}\leanok
[Pole-free at a prime divisor gives bounded adic valuation]
  \label{lem:chart_valuation_le_one_of_order_nonneg}
  \lean{AlgebraicGeometry.chart_valuation_le_one_of_order_nonneg}
  \uses{def:order_at_point, lem:stalk_isDVR_of_smooth}
  Let \(V\) be an affine integral scheme that is regular in codimension one, and
  write \(A = \Gamma(V, \top)\) for its coordinate ring and \(K = K(V) =
  \Frac(A)\) for its function field. Let \(Z\) be a prime divisor of \(V\),
  i.e.\ a codimension-one point, and let \(\mathfrak p \subseteq A\) be the
  corresponding height-one prime, so that the localisation \(A_{\mathfrak p}\)
  is the stalk \(\struct{V, Z}\), a discrete valuation ring
  (\cref{lem:stalk_isDVR_of_smooth}). Then for every \(g \in K\),
  \[
    \ord_Z(g) \ge 0 \;\Longrightarrow\;
    v_{\mathfrak p}(g) \le 1,
  \]
  where \(v_{\mathfrak p} = \texttt{HeightOneSpectrum.valuation}\) is the
  \(\mathfrak p\)-adic valuation of the Dedekind domain \(A\).
\end{lemma}

\begin{proof}
  \uses{def:order_at_point, lem:stalk_isDVR_of_smooth}
  \leanok
  By definition (\cref{def:order_at_point}), \(\ord_Z(g) \ge 0\) if and only
  if \(\operatorname{ordFrac}_{A_{\mathfrak p}}(g) \ge 1\) in the discrete
  valuation ring \(\struct{V, Z} = A_{\mathfrak p}\). On a discrete valuation
  ring the fractional order and the adic valuation are reciprocal,
  \(\operatorname{ordFrac} = (-)^{-1} \circ v\) in the value group; hence
  \(\operatorname{ordFrac}(g) \ge 1\) if and only if \(v(g) \le 1\) for the
  valuation of the localisation \(A_{\mathfrak p}\). Finally the
  \(\mathfrak p\)-adic valuation \(v_{\mathfrak p}\) of the Dedekind domain
  \(A\) at the height-one prime \(\mathfrak p\) agrees with the valuation of its
  localisation \(A_{\mathfrak p}\), so \(v_{\mathfrak p}(g) \le 1\).
\end{proof}

\begin{theorem}\leanok
[Algebraic Hartogs on a Dedekind chart]
  \label{thm:chart_mem_range_of_forall_order_nonneg}
  \lean{AlgebraicGeometry.chart_mem_range_of_forall_order_nonneg}
  \uses{thm:chart_isDedekindDomain, lem:mem_integers_of_valuation_le_one,
        def:order_at_point, lem:stalk_isDVR_of_smooth,
        lem:chart_valuation_le_one_of_order_nonneg}
  % SOURCE: [Hartshorne], II.6, Proposition 6.3A, p.~132 (read from
  % references/hartshorne-algebraic-geometry.pdf, PDF page 149 = doc p.132;
  % scanned-image page, transcribed character-by-character from the rendered
  % page).
  % SOURCE QUOTE: "Proposition 6.3A. Let $A$ be an integrally closed noetherian
  % domain. Then
  % $A = \bigcap_{\mathrm{ht}\,\mathfrak p = 1} A_{\mathfrak p}$
  % where the intersection is taken over all prime ideals of height 1.
  % PROOF. Matsumura [2, Th. 38, p. 124]."
  \textit{Source: Hartshorne, II.6, Proposition 6.3A, p.~132.}
  Let \(V\) be an affine, integral, locally Noetherian scheme that is regular in
  codimension one and satisfies \(\krullDim V \le 1\) (\(\texttt{Order.krullDim}\
  V \le 1\)). Write \(A = \Gamma(V, \top)\) for its coordinate ring and
  \(K = K(V) = \Frac(A)\) for its function field. Let \(g \in K\). If
  \[
    \ord_Z(g) \ge 0 \qquad \text{at every prime divisor } Z \text{ of } V,
  \]
  then \(g\) lies in the image of the structure map
  \(\Gamma(V, \top) \to K(V)\); that is, \(g\) is a global regular function,
  \(g \in (\operatorname{algebraMap} A\,K).\mathrm{range}\).
\end{theorem}

% SOURCE QUOTE PROOF: "PROOF. Matsumura [2, Th. 38, p. 124]."
\begin{proof}
  \uses{thm:chart_isDedekindDomain, lem:mem_integers_of_valuation_le_one,
        lem:chart_valuation_le_one_of_order_nonneg, def:order_at_point}
        \leanok
  By \cref{thm:chart_isDedekindDomain} the coordinate ring \(A = \Gamma(V,
  \top)\) is a Dedekind domain whose fraction field is the function field
  \(K = K(V)\). The algebraic-Hartogs anchor
  \cref{lem:mem_integers_of_valuation_le_one} then reduces the claim to a bound
  on adic valuations: if \(v_{\mathfrak p}(g) \le 1\) at every height-one prime
  \(\mathfrak p\) of \(A\), then \(g \in (\operatorname{algebraMap} A\,K).\mathrm{range}\).
  So it suffices to verify \(v_{\mathfrak p}(g) \le 1\) for every such
  \(\mathfrak p\).

  Fix a height-one prime \(\mathfrak p\) of \(A\). Under the homeomorphism
  \(V \cong \Spec A\) the height-one primes of \(A\) correspond exactly to the
  codimension-one points of \(V\), i.e.\ to the prime divisors \(Z\) of \(V\)
  (the same correspondence underlying the finiteness of the order support on an
  affine chart): the localisation \(A_{\mathfrak p}\) is the stalk
  \(\struct{V, Z}\), a discrete valuation ring by regularity in codimension one
  (\cref{lem:stalk_isDVR_of_smooth}). The hypothesis gives \(\ord_Z(g) \ge 0\)
  at this prime divisor, and the per-prime translation
  \cref{lem:chart_valuation_le_one_of_order_nonneg} converts it to
  \(v_{\mathfrak p}(g) \le 1\): the divisorial order is the logarithm of the
  fractional order in \(\struct{V, Z}\) (\cref{def:order_at_point}), and on a
  discrete valuation ring \(\operatorname{ordFrac} = (-)^{-1} \circ v\), so
  \(\ord_Z(g) \ge 0 \Leftrightarrow \operatorname{ordFrac}(g) \ge 1
  \Leftrightarrow v(g) \le 1\), the adic valuation of \(A\) at \(\mathfrak p\)
  agreeing with that of its localisation \(A_{\mathfrak p}\).

  As \(\mathfrak p\) was an arbitrary height-one prime, \(v_{\mathfrak p}(g) \le
  1\) for all of them; \cref{lem:mem_integers_of_valuation_le_one} then yields
  \(g \in (\operatorname{algebraMap} A\,K).\mathrm{range}\), i.e.\ \(g\) is a
  global regular function on \(V\).
\end{proof}
